\section{Introduction}

Interest in understanding how traits change dynamically within populations has driven a substantial body of research examining the role of stochasticity and noise in nature. Every biological population is a system far from equilibrium \footnote{It means that a trait distribution is outside an attractor or at a fixed point, but these systems change due to directional selection, transients, or non-stationary dynamics (such as cycles, stochasticity, or frequency dependence).}, where genetic information and traits flow, interact, and reorganize across time and space. This continuous interplay and variation is what we call evolution. Evolution requires reproduction of information (\cite{lamarckPhilosophieZoologique1809} as discussed in Nowak \cite{nowakEvolutionaryDynamicsExploring2006}), an imperfect process in which error and mutations can occur, these mutations act as a source of stochasticity that perturbs the dynamics of the population  as is discussed in Nowak chapter 2 \cite{nowakEvolutionaryDynamicsExploring2006}. Dynamical systems and stochastic processes offer a framework for thinking about population evolution as a complex system.  In this context, statistical physics provides a formal language for stochasticity and collective change, linking evolutionary dynamics to statistical-style analogies.\footnote{As a curious aside, Barton and Coe \cite{bartonApplicationStatisticalPhysics2009} develop a formal—not literal—analogy between Wright’s stationary allele-frequency distribution and the Boltzmann distribution: population size plays the role of inverse temperature, and log mean fitness that of a negative energy.} In this sense, biophysics is not merely an eclectic combination of disciplines; rather, it is defined by the questions it seeks to answer and its pursuit of general physical principles that transcend the specific details of individual biological systems \cite{bialek2012biophysics}. \\

Yet biology offers numerous cases that challenge this classical expectation. There exist phenotypes that are disadvantageous with predation and/or costly to maintain across generations. Despite this fact, these specific phenotypes persist in the population over generations. Examples include brightness of fireflies \textit{Photinus collustrans}, the long tails of male hummingbirds \textit{Spathura underwoodi}, the scents left by some species of mice \textit{Peromyscus manicultatus bairdi}, and even the chuck sound of certain frogs \textit{Gryllodes integers}\footnote{Refer to the table \ref{tab:1} on the appendix for other examples.}. Each of these species has its specific predator  (\cite{ryan1985,  cushing1985estrous} cited by \cite{magnhagen1991}). This is not consistent with classical evolutionary beliefs, where it is expect that non-adaptive phenotypes would disappear from the population. This can be explained by the deed that these genetic traits are sexually selected, meaning that individuals with these traits have a reproductive advantage, allowing these phenotypes to persist in the population despite associated costs. This data provoked Darwin to develop his theory of sexual selection talked in \cite{andersson1994} (hay que revisar que parte del libro habla de darwin y su teoria de la seleccion natural como tal). \\

A central complication is that the reproductive advantage of a trait often depends on what other individuals among population do. This is the domain of evolutionary game theory \textbf{(EGT)}, which models the fitness of a strategy as a function of its frequency and of the strategies adopted by the rest of the population \cite{nowakEvolutionaryDynamicsExploring2006}. Frequency-dependent interactions can produce stable coexistence of phenotypes, alternative equilibria, or the dominance of a single strategy, none of which can be captured by a fitness that is fixed in advance. Fitness landscapes complement this view, by representing fitness as a function over the space of phenotypes, they provide a geometric intuition for how populations move toward stable points, coexist, or become dominated by a particular phenotype, while also accommodating stochastic phenomena such as mutation inherent to reproduction \cite{bartonApplicationStatisticalPhysics2009}.\\


\subsection{Objetives}
\subsubsection{General Objective}

To develop a theoretical and computational model to investigate the evolution, stability, and coexistence of phenotypes in populations subject to strong sexual selection and predation. The model integrates evolutionary game theory, replicator dynamics, and the stochastic Moran process to characterize evolutionary trajectories through a fitness landscape incorporating demographic noise. The framework is applied to male guppies (\textit{Poecilia reticulata}), using empirical evidence to inform its biological assumptions and parameters.

\subsubsection{Specific Objectives}
\begin{enumerate}
    \item Define the phenotypes, fitness functions, and biological and behavioural hypotheses for the guppy population, and build a payoff matrix that captures the interactions between two phenotypes under different environmental conditions (presence and absence of predation).
    
    \item  Implement the Moran process among sexual selected population as an initial test of the theoretical hypotheses, including a graphical representation of the growth of the preferred sexual trait across successive generations under predation.
    
    \item Determine the conditions that allow for phenotype coexistence, considering the effects of predation and sexual selection.
    
    \item  Explore environmental factors to outline fitness landscapes that can be incorporated into the model, providing a more comprehensive understanding of the dynamics of sexual selection for real populations.
\end{enumerate}


This approach integrates evolutionary game theory, the Moran process under weak selection, and the fitness landscape to investigate the persistence of costly sexual traits \cite{nowakEvolutionaryDynamicsExploring2006,andersson1994}, bridging classical sexual selection models and stochastic dynamics. Applied to guppies (\textit{Poecilia reticulata}), it provides insight into how predation influences phenotype coexistence \cite{houdeMATECHOICEBASED1987,magnhagen1991,houdeSexColorMate1997,endlerNATURALSELECTIONCOLOR1980}, while offering a reproducible framework adaptable to other species. We hypothesize that, under weak selection, two male phenotypes can coexist stably when the reproductive benefits of sexual attractiveness offset predation costs, whereas the preferred phenotype is expected to dominate in the absence of predation.\\


To conclude this section, the document is structured as follows: Chapter 3 presents the theoretical framework, including a review of relevant literature and the development of the model. Chapter 4 describes the methodology used to implement the model and conduct simulations. Chapter 5 presents the results obtained from the simulations and mathematical results, along with an pre-analysis of their implications. Chapter 6 discusses the findings in relation to existing literature and outlines potential avenues for future research. Finally, Chapter 8 provides an appendix with supplementary information and data supporting the study's conclusions.\\